\section{Related Work}
\label{sec:related}

\subsection{Challenges in Agentic Systems}
\label{sec:related_challenges_in_mas}
The promising capabilities of agentic systems have inspired research into solving specific challenges. For instance, Agent Workflow Memory \cite{wang2024agentworkflowmemory} addresses long-horizon web navigation by introducing workflow memory. DSPy \cite{khattab2023dspycompilingdeclarativelanguage} tackles issues in programming agentic flows, while StateFlow \cite{wu2024stateflowenhancingllmtasksolving} focuses on state control within agentic workflows to improve task-solving capabilities. Several surveys also highlight challenges and potential risks specifically within MAS \cite{han2024llmmultiagentsystemschallenges, hammond2025multiagentrisksadvancedai}. While these works meaningfully contribute towards understanding specific issues or providing high-level overviews, they do not offer a fine-grained, empirically grounded taxonomy of \textit{why} MAS fail across diverse systems and tasks. Numerous benchmarks also exist to evaluate agentic systems \cite{jimenez2024swebench, peng2024survey, wang2024battleagentbench, anne2024harnessing, bettini2024benchmarl, long2024teamcraft}. These evaluations are crucial but primarily facilitate a top-down perspective, focusing on aggregate performance or high-level objectives like trustworthiness and security \cite{liu2023trustworthy, yao2024survey}. Our work complements these efforts by providing a bottom-up analysis focused on identifying specific failure modes in MAS.

\subsection{Design Principles for Agentic Systems}
Several works highlight challenges in building robust agentic systems and suggest design principles, often focused on single-agent settings. For instance, Anthropic's blog post emphasizes modular components and avoiding overly complex frameworks \cite{Anthropic_2024}. Similarly, \citet{kapoor2024aiagentsmatter} demonstrates how complexity can hinder practical adoption. Our work extends these insights to the multi-agent context. By systematically collecting and analyzing a large corpus of MAS failure instances within \dataset{}, and by developing \taxonomy{}, we provide not only a structured understanding of \textit{why} MAS fail but also empirical data from \dataset{} to support the development and validation of more robust design principles for MAS. This aligns with the call for clearer specifications and design principles \cite{stoica2024specifications}.

\subsection{Related Datasets and Taxonomy}
Despite the growing interest in LLM agents, dedicated research systematically characterizing their failure modes remains limited, particularly for MAS. While \citet{bansal2024challenges} catalogs challenges in human-agent interaction, our contribution focuses specifically on failures within autonomous MAS execution. Other related work includes taxonomies for evaluating multi-turn LLM conversations \cite{Bai_2024} or specific capabilities like code generation \cite{SongD_code_tax}. These differ significantly from our goal of developing a generalizable failure taxonomy for multi-agent interactions and coordination. 

Further related efforts aim to improve MAS through different approaches. AgentEval \cite{arabzadeh2024assessingverifyingtaskutility} proposes a framework using LLM agents to define and quantify multi-dimensional evaluation criteria reflecting task utility for end-users. AGDebugger \cite{epperson2025interactive} introduces an interactive tool enabling developers to debug and steer agent teams by inspecting and editing message histories. And currentwork by \citet{zhang2025agentcausestaskfailures} present the Who\&When dataset and MAS debugger, which focuses on summarizing failures for specific task items by attributing them to particular agents and error steps. 

Thus, \dataset{} and \taxonomy{} represent, to our knowledge, the first empirically derived, comprehensive dataset and taxonomy focused specifically on MAS failures focus on failure patterns. Identifying these patterns highlights the need for continued research into robust evaluation metrics and mitigation strategies tailored for the unique challenges of MAS.